\chapter{Host Company}
\label{chap:host-company}

\section{Company Presentation}

This section presents ST2i itself: its history, business sector, and
internal organization.

\subsection{Company Overview}

ST2i (Soci\'et\'e Tunisienne d'Informatique pour l'Ing\'enierie) is a
Tunisian information technology company founded in 1992 and a subsidiary of
the STUDI Group, an international engineering group established in 1970
with operations across Africa and Europe. With a multidisciplinary team of
engineers, consultants, and technical experts, ST2i delivers digital
transformation projects for both public and private organizations in
Tunisia and internationally.

\begin{figure}[H]
  \centering
  \includegraphics[width=0.7\textwidth]{st2i-office.png}
  \caption[ST2i offices]{ST2i's offices.}
  \label{fig:st2i-office}
\end{figure}

\subsection{Business Sector and Organization}

ST2i operates in the Information and Communication Technology (ICT) sector
as an IT services and consulting company, spanning information systems
consulting, custom software development, system integration, Geographic
Information Systems (GIS), digital transformation consulting,
telecommunications solutions, and technical support and training, across
industries including public administration, telecommunications, banking,
healthcare, education, energy, transportation, and infrastructure. Its
organizational structure brings together consulting, functional, and
technical departments (Executive Management, Business Development, Project
Management, Consulting and Business Analysis, Software Development,
Infrastructure and Systems Integration, Quality Assurance and Technical
Support), with project teams typically pairing a Project Manager and
Business Analysts with Software Engineers, QA Engineers, and Technical
Experts.

\section{Existing Environment}

ST2i bids on public and international tenders (often funded by institutions
such as the World Bank or GIZ) that specify one or more required expert
positions. Matching internal experts to these positions is currently done
manually: a business analyst reads the Terms of Reference (TDR) and searches
the internal CV database by memory or keyword search, with no systematic
scoring, no anomaly checking, and no audit trail explaining why a given
expert was or was not selected. Real TDRs are often 100 to 200+ pages; the
analyst currently relies on a paid ChatGPT Plus subscription just to be able
to upload and read documents that size, an ongoing licensing cost incurred
to compensate for a missing internal tool, not a real solution to the
staffing problem itself.

\subsection{Identified Limitations}

Table~\ref{tab:limitations} summarizes the five recurring limitations of
this manual process, each with a direct operational or financial cost.

\begin{table}[H]
  \centering
  \caption{Limitations of the current manual matching process}
  \label{tab:limitations}
  \begin{tabular}{@{}p{3.4cm}p{14.3cm}@{}}
    \toprule
    Slow and costly & Manually reading a TDR and cross-referencing dozens of CVs does not scale with the number of active tenders; the ChatGPT Plus workaround for the largest documents is a recurring cost for a problem an internal tool should solve outright. \\
    Imperfect memory & No analyst holds an accurate, up-to-date record of every employee's diplomas, certifications, and project history. People routinely misremember details of their own CV, let alone a colleague's. \\
    Inconsistent & Two analysts reviewing the same TDR may reach different conclusions, with no documented rationale to reconcile them. \\
    Error-prone at scale & The larger the TDR and the more personnel to check against, the more likely a genuine match or a disqualifying gap gets missed: a missed qualified expert can mean passing on a winnable bid, a missed gap can mean committing to unqualified staff after the bid is won. \\
    No institutional memory & Past staffing decisions are not captured in a form that could calibrate future ones, so the same manual evaluation starts from zero every time. \\
    \bottomrule
  \end{tabular}
\end{table}

Taken together, these limitations cap how many tenders ST2i can safely
pursue at once, and translate directly into missed or mishandled bids: lost
revenue.

\section{Problem Statement}

Given a TDR describing one or more required expert positions and ST2i's
internal CV database, how can the company reliably determine, for each
position, which employees are qualified, with enough explanation that a
human reviewer can trust and, when necessary, override the result, at a
speed and consistency that does not degrade as the company and its tender
volume grow?

\section{Proposed Solution}

AdéquIA automates this matching step end to end: it parses CVs and TDRs
into structured data regardless of document length, retrieves and scores
candidates against each required position using a combination of
rule-based hard filters, weighted embedding similarity, and LLM judgment
for qualitative criteria, flags anomalies (date inconsistencies, expired
certifications, implausible claims), and produces a citation-backed
rationale for every score. A human reviewer stays in the loop at every
decision point: every automated ranking can be approved, rejected, or
overridden, and that decision is written back to inform future scoring.
Where the manual process degrades as ST2i grows, this one does not: the
cost of evaluating one more candidate or one more position is a few
seconds of automated computation, not a proportionally larger share of an
analyst's attention. The remainder of this report focuses on how it was
built, sprint by sprint.
